% =====================================================================
%  Conclusion
%  -------------------------------------------------------------------
%  Summarises validation, benchmark outcome, and limitations; outlines
%  plausible future directions.
% =====================================================================

\section{Conclusion}\label{sec:conclusion}

Blaze2D is a from-scratch two-dimensional plane-wave eigensolver that
sets out to deliver MPB-grade accuracy with a workflow tuned for the
high-throughput regime: many crystal configurations, moderate plane-wave
cutoffs, and direct access to Bloch coefficients from a Python frontend.
The construction proceeds along three areas.

\emph{Accuracy}: analytic subpixel smoothing of the
permittivity reproduces MPB's $\epsilon_{\bm{G},\bm{G}'}$ to within
$\sim 10^{-4}$ in the mean even at the coarsest grids
(\cref{fig:material_map_comparison}); eigenvalues then agree to
$\sim 2 \times 10^{-5}$ at the resolutions relevant for production
runs (\cref{fig:eigenvalue_comparison}); and both solvers exhibit the
same convergence rate as the plane-wave cutoff is refined
(\cref{fig:convergence_study}). The remaining discrepancies are
localised at degenerate eigenvalues, where any orthonormal basis of
the degenerate subspace is equally valid, or at dielectric interfaces,
where the two solvers discretise the boundary slightly differently.
Neither family is a numerical disagreement; both are explicable from
solver conventions alone.

\emph{Throughput}: the central performance lever is not
algorithmic but architectural. Profiling reveals that LOBPCG on PWE
operators is memory-bandwidth-bound, not compute-bound, so a
mixed-precision storage strategy yields a near-clean $2\times$
wall-clock speedup by storing bulk arrays in single precision and
promoting only numerically sensitive reductions to double precision,
without altering the eigenvalue trajectory
(\cref{fig:blaze2d_benchmarks,fig:blaze2d_iterations_memory}). Rust's
type-safe concurrency keeps multi-core scaling predictable inside
parameter sweeps, where solves are independent and parallel.

\emph{Accessibility}: implementing the TE and TM operators
directly, rather than projecting them out of a general
curl--curl operator, keeps the per-polarization data path short
and makes operator-level matrix elements available to the Python layer
through Rust's foreign-function interface. The same operator objects
that drive the eigensolve are reused downstream, so derived quantities
(group velocities, Berry connections, effective-mass tensors) are
computed against exactly the discretization the eigenvalues were
obtained on.

Several limitations should be acknowledged. Blaze2D is two-dimensional
by design and currently exposes only the in-plane TE/TM
decomposition; fully three-dimensional or out-of-plane band
calculations require the more general curl--curl formulation and are
left to upstream tools. Band tracking across avoided crossings is not
attempted: the solver reports the $N$ lowest eigenvalues at each
$\bm{k}$-point without enforcing adiabatic continuity, which simplifies
the inner loop but pushes the responsibility for band labelling onto
the analysis stage. Small but visible asymmetries in degenerate
mode profiles (\cref{fig:u2_comparison}) have not been quantified
here; they are expected to be artefacts of the chosen basis within the
degenerate subspace, but a rigorous gauge-fixing diagnostic remains
future work.

Natural extensions follow from each of these limits. A
three-dimensional curl--curl variant, an explicit band-tracking layer
based on overlap matching across neighbouring $\bm{k}$-points, and a
GPU back end that exploits the same mixed-precision idea on
single-precision-favourable hardware are all plausible next steps. On
the application side, exposing Bloch-level data through a stable
Python ABI opens the door to coupling Blaze2D with higher-level
photonic-design frameworks, including inverse design, perturbative
envelope theories, and topological-invariant computations, in a way that
treats the solver as a service rather than a monolith.

Within its stated scope of two-dimensional photonic crystals at
moderate resolution, swept over many configurations, with full
Bloch-level data exposed, Blaze2D matches MPB on accuracy, exceeds
it on throughput in the mixed-precision regime, and provides a
cleaner programming interface for downstream analyses. The solver is
released under an open-source licence with the intent that other
groups can re-use it as a drop-in component of their own pipelines.
